\exercisesection{9}

\begin{enumerate}
\def\labelenumi{\arabic{enumi}.}
\tightlist
\item
  \begin{enumerate}
  \def\labelenumii{(\alph{enumii})}
  \tightlist
  \item
    设 K 为域（不必交换），T 为 K 上与环结构相容的局部紧拓扑；证明 T 与 K 的域结构相容。（使用 R. Ellis 定理（一般拓扑，第 X 章 § 3，习题 25），或直接仿照积分论第七章 § 1，第 10 节及本段命题 1 的证明。）
  \end{enumerate}
\end{enumerate}

\begin{enumerate}
\def\labelenumi{(\alph{enumi})}
\setcounter{enumi}{1}
\tightlist
\item
  给出交换域 K 上的一个局部紧拓扑的例子，它与 K 的加法群结构相容，但不与其环结构相容。（取 K 为紧整环 A 的分式域（例如 A 为形式幂级数环 \(k[[X, Y]]\)，其中 k 有限），并取 K 中 0 的邻域在 A 中的 0 的邻域作为基本邻域系。）
\end{enumerate}

¶ 2. (a) 在空间 \(\mathbf{R}^n\)（\(n \geqslant 2\)）中，令 U 为补集非空的非空开集；证明 U 的边界含有非空完美集（参见一般拓扑，第一章 § 9，习题 17）。

\begin{enumerate}
\def\labelenumi{(\alph{enumi})}
\setcounter{enumi}{1}
\item
  由 (a) 推出：若 \(\mathbf{A}\) 是 \(\mathbf{R}^n\) 的处处稠密子集，并且与 \(\mathbf{R}^n\) 中每个完美集相交，则 \(\mathbf{A}\) 是凸的。
\item
  证明  中存在一个处处稠密的子域 \(\mathbf{C}\)，它连通且局部连通，并且是 \(\mathbf{K}\) 的纯超越扩张（应用一般拓扑第九章 § 5 的习题 18 (b)、(c)，利用 (b) 以及集合论第三章 § 6，习题 24 中描述的方法，通过超限归纳构造 \(\mathbf{Q}\)）。推出 \(\mathbf{K}\) 中存在包含  的子域 ，它（代数地）同构于 \(\mathbf{K}' \supset \mathbf{K}\)，且连通并局部连通。\(\mathbf{C}\)\(\mathbf{R}\)，且连通并局部连通。
\item
  利用 (c) 证明 C 上存在一种连通且局部连通空间的拓扑，它与域结构相容，并且在此拓扑下 C 的完备化是 C 上的一个代数，且是两个同构于 C 的域的直复合。
\end{enumerate}

\begin{enumerate}
\def\labelenumi{\arabic{enumi}.}
\setcounter{enumi}{2}
\tightlist
\item
  设 K 为全不连通的非离散局部紧交换域；令 A 为 K 上模 \(_{K}\) 的绝对值的环，令 U 为 A 的单位群。对每个整数 n \textgreater{} 0，令 \(_{n}U\) 表示 K 中 n 次单位根所成的群，令 \(U^{n}\) 表示 U 的由 U 中元素的 n 次幂组成的子群。证明若 n 与 K 的特征互素，则 \(U^{n}\) 是 U 的开子群，并且
\end{enumerate}

\[
\operatorname{Card} (\mathrm{U} / \mathrm{U} ^ {n}) = \operatorname{mod} _ {\mathrm{K}} (n). \operatorname{Card} (_ {n} \mathrm{U})
\]

（使用积分论第七章 §2 的习题 14 以及交换代数第三章 §4，第 6 节，定理 2 的推论 1，证明若 m 是 A 的极大理想，则  在  下的像为 \(x \mapsto x^n\)，其中 k 充分大。）\(1 + \mathfrak{m}^k\)\(1 + n \cdot \mathfrak{m}^k\)\(k\)，其中 k 充分大。）

\begin{enumerate}
\def\labelenumi{\arabic{enumi}.}
\setcounter{enumi}{3}
\tightlist
\item
  \begin{enumerate}
  \def\labelenumii{(\alph{enumii})}
  \tightlist
  \item
    设 K 为交换域，v 为 K 上的赋值，且 K 关于 v 为 Hensel（§ 8，习题 6）。进一步假设 K 与 v 的剩余域 k 的特征均为 0。证明 v 的环 A 中存在子域 \(K_{0}\)，使得典范映射 \(A \rightarrow k\) 限制在 \(K_{0}\) 上时，是从 \(K_{0}\) 到 k 的同构。（令 H 为 A 的子域，使 H 在典范映射下的像同构于 k 的子域 E；证明若 \(E \neq k\)，则 A 中存在 \(\alpha \notin H\)，使得 K 的子域 \(\mathrm{H}(\alpha)\) 包含于 A，且典范同构于 \(\mathrm{E}(\bar{\alpha})\)，其中 \(\bar{\alpha}\) 是 \(\alpha\) 在 k 中的类；根据 \(\bar{\alpha}\) 在 E 上代数或超越分两种情形。）
  \end{enumerate}
\end{enumerate}

\begin{enumerate}
\def\labelenumi{(\alph{enumi})}
\setcounter{enumi}{1}
\tightlist
\item
  进一步假设 v 是离散赋值且 K 关于 v 完备；由 (a) 推出 K 同构于形式幂级数域 \(k((\mathbf{T}))\)。
\end{enumerate}

\begin{enumerate}
\def\labelenumi{\arabic{enumi}.}
\setcounter{enumi}{4}
\tightlist
\item
  \begin{enumerate}
  \def\labelenumii{(\alph{enumii})}
  \tightlist
  \item
    设 K 为交换域，v 为 K 上的高度 1 赋值，且 K 关于 v 完备，A 为 v 的环；进一步假设 v 的剩余域 k 是完美域且特征为 p \textgreater{} 0。对每个元素  及每个整数 n，令 \(\xi \in k^{*}\) 为 A 中属于类 \(x_{n}\) 的元素；证明序列 \(\xi^{p-n}\) 是 A 中的 Cauchy 序列，其极限与在类 \((x_{n}^{p^{n}})\) 中选择  的方式无关。若该极限记作 \(x_{n}\)，证明 \(\xi^{p-n}\)\(\phi(\xi)\) 是从乘法群 \(\phi\) 到乘法群 \(k^{*}\) 的唯一同构 u，使得对所有 \(K^{*}\)，\(\xi \in k^{*}\) 是 A 中属于类 \(u(\xi)\) 的元素。\(\xi\) 的元素。
  \end{enumerate}
\end{enumerate}

\begin{enumerate}
\def\labelenumi{(\alph{enumi})}
\setcounter{enumi}{1}
\item
  若  的特征也为 p，证明将 \(\mathbf{K}\) 通过 \(p\)\(\phi\) 延拓到 k 后，它是从域 \(k\)\(\phi(0) = 0\) 到 \(k\) 的一个子域的同构。推出第 3 节定理 1 (iii) 的一个新证明。\(\mathbf{K}\) 的一个子域的同构。推出第 3 节定理 1 (iii) 的一个新证明。
\item
  假设  有限。证明若 \(k\) 与 p 互素，则 \(r\) 中元素的 r 次幂所成群  在 \(p\)\((\mathbf{K}^*)^r\) 中具有有限指数（使用 Hensel 引理）。进一步若 v 离散且 \(r\)\(\mathbf{K}^*\) 的特征为 0，则证明无论 \(\mathbf{K}^*\)\(v\)\(\mathbf{K}\) 如何同样成立（注意 \(r\) 的每个元素都是 p 次幂）。\(1 + p^2\mathbf{A}\)\(p\) 的每个元素都是 p 次幂）。
\end{enumerate}

